\section{结论}
本文建立了从条件溶解度建模到交互预测的流程。按药物母体划分数据，并以母体等权 MAE 选择配置，使验证目标与未见母体预测任务一致。开发集比较中，TabPFN 与丰富描述符 ExtraTrees 的双分支组合取得最低观测误差。折间权重及低溶解度端误差表明，后续验证应关注融合收益在不同母体上的稳定性，以及低溶解度样本的预测校准。
